% preamble/settings.tex — Einstellungen, Konfigurationen, Abkürzungen
% ═════════════════════════════════════════════════════════════════════

% ── Metadaten (hier anpassen!) ────────────────────────────────────────
\newcommand{\docTitle}{Fine-Tuning von NVIDIA GR00T N1.6 für humanoide Roboter-Manipulation}
\newcommand{\docSubtitle}{Block-Stacking mit dem Unitree G1 und DEX3-Hand}
\newcommand{\docAuthor}{Luca Müncke}
\newcommand{\docMatrikel}{u28320}
\newcommand{\docBetreuer}{[Name des Betreuers]}
\newcommand{\docHochschule}{[Name der Hochschule]}
\newcommand{\docFachbereich}{[Fachbereich / Institut]}
\newcommand{\docStudiengang}{[Studiengang]}
\newcommand{\docAbgabedatum}{[TT.\,MM.\,JJJJ]}

% ── Hyperref-Konfiguration ────────────────────────────────────────────
\hypersetup{
  colorlinks   = true,
  linkcolor    = NavyBlue,
  citecolor    = NavyBlue,
  urlcolor     = BurntOrange,
  filecolor    = NavyBlue,
  pdftitle     = {\docTitle},
  pdfauthor    = {\docAuthor},
  pdfsubject   = {Projektarbeit},
  pdfkeywords  = {Humanoider Roboter, GR00T, Fine-Tuning, Imitation Learning,
                  Unitree G1, DEX3},
  pdflang      = {de-DE},
  pdfcreator   = {LuaLaTeX mit hyperref},
}

% ── cleveref-Bezeichnungen (Deutsch) ─────────────────────────────────
\crefname{figure}{Abb.}{Abb.}
\Crefname{figure}{Abbildung}{Abbildungen}
\crefname{table}{Tab.}{Tab.}
\Crefname{table}{Tabelle}{Tabellen}
\crefname{equation}{Gl.}{Gl.}
\Crefname{equation}{Gleichung}{Gleichungen}
\crefname{chapter}{Kap.}{Kap.}
\Crefname{chapter}{Kapitel}{Kapitel}
\crefname{section}{Abschn.}{Abschn.}
\Crefname{section}{Abschnitt}{Abschnitte}
\crefname{listing}{Listing}{Listings}
\Crefname{listing}{Listing}{Listings}
\crefname{appendix}{Anh.}{Anh.}
\Crefname{appendix}{Anhang}{Anhänge}

% ── Caption-Stil ──────────────────────────────────────────────────────
\captionsetup{
  font      = small,
  labelfont = {bf, sf},
  labelsep  = period,
  skip      = 6pt,
}
\captionsetup[sub]{
  font      = footnotesize,
  labelfont = sf,
}

% ── minted-Standardkonfiguration ─────────────────────────────────────
\setminted{
  bgcolor    = gray!8,
  fontsize   = \small,
  linenos    = true,
  breaklines = true,
  breakanywhere,
  frame      = leftline,
  framesep   = 2mm,
  rulecolor  = NavyBlue!50,
  numbersep  = 4pt,
  tabsize    = 4,
}

% ── siunitx ──────────────────────────────────────────────────────────
\sisetup{
  output-decimal-marker = {,},
  group-separator       = {\,},
  group-minimum-digits  = 4,
}

% ── KOMA-Script Schriftsatz ───────────────────────────────────────────
\setkomafont{chapter}{\LARGE\bfseries\sffamily}
\setkomafont{section}{\large\bfseries\sffamily}
\setkomafont{subsection}{\large\sffamily}
\setkomafont{subsubsection}{\normalsize\itshape\sffamily}
\setkomafont{disposition}{\sffamily}
\addtokomafont{chapterprefix}{\color{NavyBlue!60}}
\addtokomafont{chapter}{\color{NavyBlue}}

% ── Kopf- und Fußzeile ────────────────────────────────────────────────
\pagestyle{scrheadings}
\clearpairofpagestyles
\ihead{\headmark}
\ofoot[\pagemark]{\pagemark}
\setkomafont{pageheadfoot}{\small\sffamily}
\setkomafont{pagenumber}{\sffamily}
\automark[section]{chapter}

% ── Aufzählungen ─────────────────────────────────────────────────────
\setlist{noitemsep, topsep=0.5\baselineskip}
\setlist[itemize,1]{label=\textbullet}
\setlist[itemize,2]{label=\textendash}
\setlist[enumerate,1]{label=\arabic*.}

% ── tcolorbox-Umgebungen ─────────────────────────────────────────────
\tcbuselibrary{breakable, skins}

\newtcolorbox{infobox}[1][Hinweis]{
  enhanced, breakable,
  colback  = NavyBlue!5,
  colframe = NavyBlue!60,
  fonttitle = \bfseries\sffamily,
  title    = {#1},
}
\newtcolorbox{warningbox}[1][Achtung]{
  enhanced, breakable,
  colback  = BrickRed!5,
  colframe = BrickRed!60,
  fonttitle = \bfseries\sffamily,
  title    = {#1},
}
\newtcolorbox{tipbox}[1][Tipp]{
  enhanced, breakable,
  colback  = OliveGreen!5,
  colframe = OliveGreen!60,
  fonttitle = \bfseries\sffamily,
  title    = {#1},
}

% ── acro-Konfiguration ────────────────────────────────────────────────
\acsetup{make-links = true}

\DeclareAcronym{vla}{
  short = VLA,
  long  = Vision-Language-Action,
}
\DeclareAcronym{llm}{
  short = LLM,
  long  = Large Language Model,
}
\DeclareAcronym{il}{
  short = IL,
  long  = Imitation Learning,
}
\DeclareAcronym{rl}{
  short = RL,
  long  = Reinforcement Learning,
}
\DeclareAcronym{gpu}{
  short = GPU,
  long  = Graphics Processing Unit,
}
\DeclareAcronym{hpc}{
  short = HPC,
  long  = High-Performance Computing,
}
\DeclareAcronym{slurm}{
  short = SLURM,
  long  = Simple Linux Utility for Resource Management,
}
\DeclareAcronym{api}{
  short = API,
  long  = Application Programming Interface,
}
\DeclareAcronym{vram}{
  short = VRAM,
  long  = Video Random Access Memory,
}
\DeclareAcronym{hf}{
  short = HF,
  long  = Hugging Face,
}
\DeclareAcronym{cnn}{
  short = CNN,
  long  = Convolutional Neural Network,
}
\DeclareAcronym{moe}{
  short = MoE,
  long  = Mixture of Experts,
}
\DeclareAcronym{vit}{
  short = ViT,
  long  = Vision Transformer,
}
\DeclareAcronym{dof}{
  short = DoF,
  long  = Degrees of Freedom,
}
\DeclareAcronym{fps}{
  short = FPS,
  long  = Frames per Second,
}
\DeclareAcronym{wandb}{
  short = W\&B,
  long  = Weights \& Biases,
}
